\documentclass[12pt,a4paper]{article}
\usepackage[T1]{fontenc}
\usepackage[utf8]{inputenc}
\usepackage[spanish]{babel}
\usepackage{geometry}
\usepackage{booktabs}
\usepackage{xcolor}
\usepackage{listings}
\usepackage{hyperref}
\usepackage{enumitem}
\usepackage{titlesec}
\usepackage{parskip}

\geometry{margin=2.5cm}

\hypersetup{
  colorlinks=true,
  linkcolor=black,
  urlcolor=blue!70!black,
}

\lstdefinestyle{bash}{
  language=bash,
  basicstyle=\small\ttfamily,
  backgroundcolor=\color{gray!10},
  frame=single,
  rulecolor=\color{gray!40},
  breaklines=true,
  showstringspaces=false,
}

\title{%
  \textbf{Instalación de Steam}\\[0.4em]
  \large Procedimiento y resolución de errores en\\
  ASUS ROG Zephyrus G14 GA403UV con CachyOS
}
\author{pc-config}
\date{\today}

\begin{document}

\maketitle
\tableofcontents
\newpage

% ─────────────────────────────────────────────
\section{Contexto}

El equipo es un \textbf{ASUS ROG Zephyrus G14 GA403UV} con \textbf{CachyOS}
(base Arch Linux) y el entorno \textbf{Omarchy} (Hyprland + Wayland).
Cuenta con configuración de GPU dual:

\begin{itemize}
  \item \textbf{iGPU}: AMD Radeon \emph{HawkPoint} --- driver \texttt{amdgpu},
        conectada directamente al panel.
  \item \textbf{dGPU}: NVIDIA GeForce RTX 4060 Laptop GPU (8 GB VRAM) ---
        driver \texttt{nvidia-open} 595.71.05, arquitectura PRIME Offload.
\end{itemize}

Steam estaba disponible en el repositorio \texttt{multilib} pero su instalación
directa fallaba por dos motivos independientes, descritos en las secciones
siguientes.

% ─────────────────────────────────────────────
\section{Problema 1: base de datos local desactualizada}

\subsection{Síntoma}

Al ejecutar \lstinline[style=bash]{sudo pacman -S steam}, pacman emitía
centenares de advertencias del tipo:

\begin{lstlisting}[style=bash]
advertencia: PAQUETE: no se encuentra la clave "%INSTALLED_DB%"
             en la base de datos local
\end{lstlisting}

Estas advertencias afectaban a prácticamente todos los paquetes instalados en
el sistema.

\subsection{Causa}

La base de datos local de pacman (\texttt{/var/lib/pacman/local/}) contenía
entradas con el marcador literal \texttt{\%INSTALLED\_DB\%} sin sustituir,
indicativo de una migración de formato de base de datos incompleta o de una
instalación inicial que no terminó de registrar correctamente los metadatos.

Aunque no bloqueaban la instalación, producían un output extremadamente ruidoso
que dificultaba detectar errores reales.

\subsection{Solución}

Forzar la resincronización completa de los repositorios para descartar
inconsistencias entre la base de datos remota y la local:

\begin{lstlisting}[style=bash]
sudo pacman -Syy
\end{lstlisting}

La doble \texttt{y} fuerza la re-descarga de todas las bases de datos remotas
aunque pacman las considere actualizadas.

% ─────────────────────────────────────────────
\section{Problema 2: paquete \texttt{lib32-libxss} con error 404}

\subsection{Síntoma}

Una vez resincronizada la base de datos, la instalación de Steam fallaba con:

\begin{lstlisting}[style=bash]
error: no se pudo obtener el archivo
       'lib32-libxss-1.2.4-2-x86_64.pkg.tar.zst'
       desde archlinux.cachyos.org: The requested URL returned error: 404
error: no se pudo obtener el archivo
       'lib32-libxss-1.2.4-2-x86_64.pkg.tar.zst.sig'
       desde fastly.mirror.pkgbuild.com: The requested URL returned error: 404
\end{lstlisting}

El error se reproducía en todos los mirrors disponibles.

\subsection{Causa}

La caché local de pacman contenía el archivo
\texttt{lib32-libxss-1.2.4-2-x86\_64.pkg.tar.zst} de una versión anterior.
Pacman lo consideraba ``actualizado'' (\emph{up to date}) pero el paquete real
en los repositorios ya había avanzado a la versión \textbf{1.2.5-1}, de modo
que al intentar descargar la firma \texttt{.sig} de la versión antigua obtenía
un 404 en todos los espejos.

\subsection{Solución}

Eliminar el archivo cacheado obsoleto y reinstalar el paquete para forzar la
descarga de la versión actual:

\begin{lstlisting}[style=bash]
sudo rm -f /var/cache/pacman/pkg/lib32-libxss*
sudo pacman -S lib32-libxss
\end{lstlisting}

El repositorio multilib proveyó correctamente la versión 1.2.5-1.

% ─────────────────────────────────────────────
\section{Procedimiento completo de instalación}

Una vez identificados y resueltos ambos problemas, la secuencia completa
para instalar Steam es la siguiente:

\begin{enumerate}
  \item Resincronizar forzosamente los repositorios:
\begin{lstlisting}[style=bash]
sudo pacman -Syy
\end{lstlisting}

  \item Limpiar la caché del paquete problemático:
\begin{lstlisting}[style=bash]
sudo rm -f /var/cache/pacman/pkg/lib32-libxss*
\end{lstlisting}

  \item Instalar Steam (acepta con \texttt{S} cuando se pregunte):
\begin{lstlisting}[style=bash]
sudo pacman -S steam
\end{lstlisting}
\end{enumerate}

Steam instala aproximadamente \textbf{48 MiB} de paquetes, incluyendo las
dependencias \texttt{lib32} necesarias para ejecutar juegos de 32 bits.

% ─────────────────────────────────────────────
\section{Uso de Steam con la GPU NVIDIA}

Por defecto Steam usa la iGPU AMD para renderizar. Para forzar la RTX 4060
en un juego concreto, establecer las opciones de lanzamiento en
\emph{Propiedades del juego}:

\begin{lstlisting}[style=bash]
prime-run %command%
\end{lstlisting}

Para activarlo globalmente en todos los juegos, ir a
\textit{Steam $\to$ Configuración $\to$ Compatibilidad} y añadir la misma
línea como opciones de lanzamiento predeterminadas.

\subsection{Overlay de monitorización}

Para visualizar el uso de GPU, CPU y framerate durante el juego:

\begin{lstlisting}[style=bash]
sudo pacman -S mangohud
\end{lstlisting}

Opciones de lanzamiento con overlay:

\begin{lstlisting}[style=bash]
MANGOHUD=1 prime-run %command%
\end{lstlisting}

\subsection{Nota sobre el MUX Switch}

El G14 GA403UV dispone de \emph{MUX switch}, lo que permite conectar la
NVIDIA directamente al panel eliminando el paso por la iGPU
(\emph{Ultimate Mode} en Windows). En Linux se gestiona con
\texttt{supergfxctl}:

\begin{lstlisting}[style=bash]
sudo pacman -S supergfxctl
sudo systemctl enable --now supergfxd
supergfxctl --mode AsusMuxDgpu   # NVIDIA directo al panel (requiere reinicio)
supergfxctl --mode Hybrid        # PRIME normal (ahorro de bateria)
\end{lstlisting}

\end{document}
